\documentclass[12pt]{article}
\usepackage[a4paper,margin=2.35cm]{geometry}
\usepackage{amsmath,amssymb,amsfonts,bm,booktabs,enumitem,graphicx,microtype,hyperref}
\hypersetup{colorlinks=true,linkcolor=blue,urlcolor=blue}
\setlength{\parskip}{0.72em}
\setlength{\parindent}{0pt}
\setlist[itemize]{topsep=0.2em,itemsep=0.15em}
\setlist[enumerate]{topsep=0.2em,itemsep=0.15em}

% Notation follows the Particle Cosmology lecture notes.
\newcommand{\be}{\begin{eqnarray}}
\newcommand{\non}{\nonumber\\}
\newcommand{\ee}{\end{eqnarray}}
\newcommand{\Hs}{\mathcal H}
\newcommand{\Rs}{\mathcal R}
\newcommand{\bk}{\mathbf k}
\newcommand{\bx}{\mathbf x}
\newcommand{\bn}{\hat{\mathbf n}}
\newcommand{\bv}{\mathbf v}
\newcommand{\dd}{\mathrm d}
\newcommand{\Mpl}{M_{\rm Pl}}
\newcommand{\keyresult}[1]{\begin{center}\fbox{\parbox{0.91\linewidth}{\textbf{Key result.} #1}}\end{center}}

\title{The Cosmic Microwave Background I\\
The expanding Universe, recombination, and acoustic oscillations}
\author{RTG 3186 Kickoff Fall Workshop, Staufen}
\date{29 September--1 October 2026}

\begin{document}
\maketitle

\begin{center}
\fbox{\parbox{0.91\linewidth}{
\textbf{Starting point and goal.}
No previous cosmology is assumed. We first build the homogeneous expanding Universe,
then follow its cooling until neutral atoms form. Finally we add small perturbations
and derive why the photon--baryon plasma oscillates. Lecture II will project those
oscillations onto the sky, derive $C_\ell$, and connect theory to real data.}}
\end{center}

Unless units are displayed explicitly, we use natural units
$c=\hbar=k_B=1$. We restore $c$ occasionally when discussing distances or emphasizing
their physical interpretation. This is the same convention used in the Particle
Cosmology notes.

% ============================================================
\section{What are we trying to explain?}

The cosmic microwave background (CMB) is electromagnetic radiation arriving from
every direction. Its mean spectrum is extremely close to a blackbody,
\be
B_\nu(T)=\frac{2h\nu^3}{c^2}\frac{1}{\exp(h\nu/k_BT)-1},
\qquad T_0\simeq 2.725\ {\rm K}.
\ee
The word \emph{microwave} describes the wavelength today. The same photons had much
higher energies in the past. The word \emph{background} means that, after removing
local emission, it fills the entire sky.

We write the temperature in direction $\bn$ as
\be
T(\bn)=\bar T\,[1+\Theta(\bn)],
\qquad \Theta(\bn)\equiv\frac{\delta T(\bn)}{\bar T}.
\ee
The sky is very uniform, but not perfectly so:
\[
|\Theta|\sim10^{-5}
\]
after removal of the dipole induced mostly by our motion. These tiny variations are
the central observable of these lectures.

\paragraph{Three facts that demand an explanation.}
\begin{enumerate}
\item Why is the spectrum blackbody, and why is its temperature only $2.725$ K today?
\item Why do photons reach us from a particular early epoch rather than from arbitrarily far back?
\item Why does the angular pattern contain a harmonic sequence of acoustic peaks?
\end{enumerate}

The answers are, respectively: an expanding thermal Universe, recombination and
decoupling, and sound waves in the primordial photon--baryon plasma.

The near-uniformity and the small variations carry complementary information. The
mean temperature and spectrum probe the thermal history, while the $10^{-5}$
anisotropies probe gravity, matter content, geometry, and the initial conditions.

\subsection{Inflation as the origin of the initial conditions}
The near-uniformity itself raises a causal question. In a decelerating hot-Big-Bang
history, many regions seen in widely separated directions on the last-scattering
surface had no time to exchange signals before recombination. Nevertheless, their
temperatures are almost identical. This is the horizon problem.

Inflation is a period of accelerated expansion,
\be
\ddot a>0,
\qquad
\frac{d}{dt}\left(\frac{1}{aH}\right)<0.
\ee
The second expression states that the comoving Hubble radius decreases. A region that
was initially small enough to be causally connected can therefore be stretched to
scales much larger than the Hubble radius. The observable Universe may then descend
from one nearly homogeneous causal patch. After inflation ends, its energy is
converted into a hot plasma and the familiar thermal evolution begins.

Inflation also supplies small deviations from homogeneity. Quantum fluctuations of
the field driving inflation are stretched to macroscopic wavelengths. In the simplest
single-field slow-roll picture,
\be
\delta\phi_k\sim\frac{H}{2\pi},
\qquad
\Rs(\bk)\simeq-\frac{H}{\dot{\bar\phi}}\delta\phi_{\bk},
\ee
where $\Rs$ is the comoving curvature perturbation. Once a wavelength is larger than
the Hubble radius, the growing adiabatic mode of $\Rs$ is approximately conserved.
It later seeds the gravitational potentials and acoustic motion studied below.

A useful parametrization of its dimensionless power spectrum is
\be
\boxed{\Delta_{\Rs}^2(k)=A_s\left(\frac{k}{k_*}\right)^{n_s-1}.}
\ee
Simple inflationary models predict perturbations that are nearly scale invariant
($n_s$ close to one), nearly Gaussian, and predominantly adiabatic. They also select
the same growing solution for every Fourier mode. This synchronizes the subsequent
acoustic phases and is essential for producing a sequence of sharp CMB peaks rather
than a washed-out spectrum. Inflation can additionally generate tensor perturbations,
whose amplitude relative to scalar perturbations is described by the tensor-to-scalar
ratio $r$.

This short discussion provides initial conditions rather than a complete derivation
of inflation. The background equations below explain what accelerated expansion
means dynamically, while Lecture II shows how $A_s$, $n_s$, and $r$ enter observable
CMB statistics.

% ============================================================
\section{The homogeneous expanding Universe}

\subsection{Homogeneity and isotropy}
On sufficiently large scales, the Universe has no preferred position and no preferred
direction. These statements form the \emph{cosmological principle}. They do not say
that the Universe is exactly smooth on small scales: stars, galaxies, and clusters are
the perturbations that we later place on top of the smooth background.

Homogeneity and isotropy constrain the spacetime metric to the Friedmann--Lema\^{i}tre--
Robertson--Walker (FLRW) form. For a spatially flat Universe,
\be
\boxed{ds^2=-dt^2+a^2(t)\delta_{ij}dx^idx^j.}
\label{eq:FLRW}
\ee
Here $t$ is cosmic time, $\bx$ is a \emph{comoving} coordinate, and $a(t)$ is the
scale factor. We set $a(t_0)=1$ today. Two observers at fixed comoving coordinates
have physical separation
\be
d_{\rm phys}(t)=a(t)|\Delta\bx|.
\ee
Expansion means that $a(t)$ grows. It is space between freely moving comoving
observers that changes; galaxies are not projectiles moving away from a central point.

\subsection{The Hubble parameter and Hubble's law}
Differentiate the physical separation:
\be
\dot d_{\rm phys}=\frac{\dot a}{a}d_{\rm phys}=H(t)d_{\rm phys},
\qquad H(t)\equiv\frac{\dot a}{a}.
\ee
In this subsection a dot means $d/dt$. For nearby objects this is Hubble's law,
$v=Hd$. The Hubble time $H^{-1}$ is the characteristic expansion timescale; the
Hubble length $cH^{-1}$ is a useful characteristic distance. It is not, in general,
the same as the age or the causal horizon.

\subsection{Matter as a perfect fluid}
The homogeneous contents are described by energy density $\rho(t)$ and pressure
$P(t)$. For a perfect fluid,
\be
T_{\mu\nu}=(\rho+P)U_\mu U_\nu+Pg_{\mu\nu}.
\ee
In the comoving frame $U^\mu=(1,0,0,0)$. Einstein's equations applied to the flat
FLRW metric give
\be
\boxed{H^2=\frac{8\pi G}{3}\rho,}
\label{eq:Friedmann}
\ee
\be
\frac{\ddot a}{a}=-\frac{4\pi G}{3}(\rho+3P).
\label{eq:acceleration}
\ee
The first Friedmann equation says that the expansion rate is determined by the total
energy density. The second says that pressure also gravitates.

It is often useful to introduce the reduced Planck mass,
\[
M_{\rm Pl}^{-2}=8\pi G,
\qquad H^2=\frac{\rho}{3M_{\rm Pl}^2}.
\]

\subsection{Local energy conservation and the equation of state}
In general relativity, energy and momentum obey the \emph{local}, covariant
conservation law $\nabla_\mu T^{\mu\nu}=0$. For a homogeneous fluid this yields
\be
\boxed{\dot\rho+3H(\rho+P)=0.}
\label{eq:continuity}
\ee
This does \emph{not} say that the total energy of an expanding Universe is constant.
A general FLRW spacetime has no time-translation symmetry that would supply a
globally conserved energy of the usual kind. In particular, the energy in a fixed
comoving region need not be constant.

The local equation has a useful thermodynamic form. A region with fixed comoving
volume has physical volume $V\propto a^3$ and fluid energy $E=\rho V$, so
\[
d(\rho V)=-P\,dV.
\]
This is the pressure--volume work relation: expansion changes the fluid energy when
$P\ne0$. For radiation, $E_\gamma\propto a^{-1}$ because photons redshift; for vacuum
energy, $E_\Lambda\propto a^3$. Neither example violates local covariant conservation.
The work analogy does not require an external reservoir into which the energy of
redshifted photons must flow.

For a separately conserved component with $P=w\rho$ and constant $w$,
Eq.~\eqref{eq:continuity} gives
\be
\boxed{\rho(a)=\rho_0a^{-3(1+w)}.}
\ee

\begin{center}
\begin{tabular}{@{}llll@{}}\toprule
Component & $w$ & Evolution & Reason\\ \midrule
Non-relativistic matter & $0$ & $\rho_m\propto a^{-3}$ & number dilution\\
Radiation & $1/3$ & $\rho_r\propto a^{-4}$ & dilution $+$ redshift\\
Vacuum energy & $-1$ & $\rho_\Lambda={\rm constant}$ & energy of empty space\\ \bottomrule
\end{tabular}
\end{center}

Inserting these scalings in the Friedmann equation gives
\be
a(t)\propto t^{1/2}\quad&\text{during radiation domination},\non
a(t)\propto t^{2/3}\quad&\text{during matter domination},\non
a(t)\propto e^{Ht}\quad&\text{during vacuum domination}.
\ee

\keyresult{The minimum background cosmology is
$ds^2=-dt^2+a^2d\bx^2$, $H=\dot a/a$,
$H^2=8\pi G\rho/3$, and $\dot\rho+3H(\rho+P)=0$. Everything in the CMB evolves on
this background.}

\subsection{Redshift and the cooling of radiation}
Consider two successive wave crests propagating through the expanding metric. Their
physical separation stretches with $a$, so
\be
\lambda\propto a,
\qquad \nu\propto a^{-1},
\qquad E_\gamma=h\nu\propto a^{-1}.
\ee
Redshift is defined by
\be
1+z\equiv\frac{\lambda_{\rm obs}}{\lambda_{\rm em}}
=\frac{a(t_0)}{a(t_{\rm em})}=\frac{1}{a(t_{\rm em})}.
\ee
A blackbody remains a blackbody under this uniform redshift, with
\be
\boxed{T(a)=\frac{T_0}{a}=T_0(1+z).}
\label{eq:Tredshift}
\ee
This connects looking to large redshift with looking back to a hotter, denser Universe.
At $z\simeq1100$, Eq.~\eqref{eq:Tredshift} gives $T\simeq3000$ K.

Radiation scales as $a^{-4}$ rather than $a^{-3}$ because there are two effects:
the number of photons per physical volume gives $a^{-3}$, while the energy of every
photon contributes one additional factor $a^{-1}$.

% ============================================================
\section{Light propagation, conformal time, and horizons}

For photon propagation it is convenient to define conformal time
\be
d\tau\equiv\frac{dt}{a(t)}.
\ee
The FLRW metric becomes
\be
ds^2=a^2(\tau)\left[-d\tau^2+\delta_{ij}dx^idx^j\right].
\ee
Apart from the overall factor $a^2$, this looks like Minkowski space. Introduce the
\emph{comoving radial coordinate} $\chi=|\bx|$, with the observer at $\chi=0$.
At a fixed time the physical radial distance is $a(\tau)\chi$; in spherical
coordinates $d\bx^2=d\chi^2+\chi^2d\Omega^2$.
A radial light ray has $d\Omega=0$ and $ds^2=0$, hence $|d\chi|=d\tau$ when $c=1$.
For a photon approaching us, $d\chi/d\tau=-1$: its radial coordinate decreases as
time increases. The positive comoving distance to an object observed at redshift
$z$ is consequently
\be
\chi(z)=\tau_0-\tau(z)=\int_0^z\frac{dz'}{H(z')}.
\label{eq:distance}
\ee

From here onward, as in the perturbation sections of the Particle Cosmology notes,
a dot denotes a conformal-time derivative. We distinguish
\be
\Hs\equiv\frac{\dot a}{a}=aH,
\ee
where $H$ is the usual cosmic-time Hubble parameter.

The comoving particle horizon at time $\tau$ is the maximum comoving distance a signal
could have travelled since the beginning,
\be
d_H^{\rm com}(\tau)=\tau-\tau_{\rm initial}\simeq\tau.
\ee
Here the last equality chooses the origin of conformal time at the beginning.
The comoving Hubble length is $(aH)^{-1}=\Hs^{-1}$. It measures an expansion
timescale expressed as a distance, whereas the particle horizon counts the full
distance light could have travelled. They are related but need not coincide.
Pressure disturbances travel more slowly than light; their corresponding causal
distance will be the \emph{sound horizon}, introduced in Section~\ref{sec:acoustics}.

\subsection{Why the CMB is a two-dimensional image}
If photons last scattered at conformal time $\tau_*$, we observe them today at
$\tau_0$ on a sphere of comoving radius
\be
\boxed{\chi_*=\tau_0-\tau_*.}
\ee
A patch of comoving transverse size $L_{\rm com}$ at last scattering appears as an
angular scale on this sphere. For a small angle,
\be
\theta\simeq\frac{L_{\rm com}}{\chi_*}.
\ee
Both lengths are comoving: using physical lengths instead multiplies numerator and
denominator by the same $a_*$. This simple geometry will later connect sound waves
in three dimensions to angular patterns on our sky.

% ============================================================
\section{A short thermal history of the photon bath}

At sufficiently high temperature the Universe is a plasma. Photons, electrons,
positrons, nuclei, neutrinos, and other species interact. A microscopic interaction
rate $\Gamma$ should be compared with the macroscopic expansion rate $H$:
\be
\Gamma\gg H &: & \text{many interactions per expansion time; equilibrium},\non
\Gamma\lesssim H &: & \text{interactions cannot keep up; decoupling}.
\ee
This single comparison connects particle physics to cosmology.

For a species in thermal equilibrium, the distribution is
\be
f(E)=\frac{1}{\exp[(E-\mu)/T]\pm1},
\ee
with $+$ for fermions and $-$ for bosons. Photons have vanishing chemical potential
in full equilibrium. The photon number and energy densities scale as
\be
n_\gamma=\frac{2\zeta(3)}{\pi^2}T^3,
\qquad \rho_\gamma=\frac{\pi^2}{15}T^4.
\ee
The observed blackbody spectrum therefore tells us that the radiation was once in
excellent thermal equilibrium. Expansion subsequently redshifted the spectrum without
changing its Planck form.

The broad sequence relevant here is
\[
\begin{split}
\text{hot ionized plasma}&\longrightarrow\text{neutral atoms form}\\
&\longrightarrow\text{photons free-stream}\longrightarrow\text{CMB observed today}.
\end{split}
\]

% ============================================================
\section{Recombination, decoupling, and last scattering}

\subsection{Why the early Universe was opaque}
Free electrons Thomson-scatter photons with cross section
\be
\sigma_T=\frac{8\pi}{3}\left(\frac{\alpha}{m_e}\right)^2
\ee
in natural units. The scattering rate is roughly
\be
\Gamma_T=n_e\sigma_T.
\ee
At early times $n_e$ is large and $\Gamma_T/H\gg1$. A photon travels only a short
mean free path $\lambda_{\rm mfp}=(n_e\sigma_T)^{-1}$ before scattering, so the plasma
is opaque. Coulomb scattering couples electrons to protons; Thomson scattering couples
electrons to photons. The three components behave approximately as one fluid.

\subsection{Deriving the Saha equation: why recombination is below 13.6 eV}
As the temperature falls, electrons and protons form neutral hydrogen,
\[
p+e^-\leftrightarrow H^0+\gamma.
\]
First neglect helium and excited hydrogen levels. Here $n_p$ counts
\emph{free protons}, $n_e$ counts \emph{free electrons}, and $n_{H^0}$ counts
neutral hydrogen atoms. A proton bound inside a neutral atom is included in
$n_{H^0}$, not in the free-proton species density $n_p$. Thus
\[
n_H=n_p+n_{H^0}
\]
is the \emph{total number density of hydrogen nuclei}, whether free or bound.
Only in a fully ionized gas is $n_H\simeq n_p$. Recombination changes how nuclei
are partitioned between free protons and neutral atoms; it does not destroy
nuclei. Their number in a comoving volume is conserved, while their physical
density $n_H$ dilutes as $a^{-3}$.

Charge neutrality gives $n_e=n_p$. The useful ionization fraction is therefore
\[
x_e\equiv\frac{n_e}{n_H},\qquad
n_e=n_p=x_en_H,\qquad n_{H^0}=(1-x_e)n_H.
\]
Normalizing instead by the free-proton density would give $n_e/n_p=1$
whenever free charges are present, and would not measure the fraction of
hydrogen that is ionized.

\paragraph{From chemical equilibrium to number densities.}
For the reaction above, equilibrium requires
\[
\mu_p+\mu_e=\mu_{H^0}+\mu_\gamma=\mu_{H^0},
\]
because a blackbody photon bath has $\mu_\gamma=0$. The atoms, protons, and electrons
are non-relativistic and dilute. Integrating their Maxwell--Boltzmann distributions
over momentum gives, with chemical potentials including rest energy,
\[
n_i=g_i\int\frac{d^3p}{(2\pi)^3}
e^{-(m_i+p^2/2m_i-\mu_i)/T}
=g_i\left(\frac{m_iT}{2\pi}\right)^{3/2}e^{(\mu_i-m_i)/T}.
\]
Substitute this expression for each of the three species before simplifying:
\[
\begin{aligned}
\frac{n_en_p}{n_{H^0}}
={}&\frac{g_eg_p}{g_{H^0}}
\left(\frac{m_em_p}{m_{H^0}}\frac{T}{2\pi}\right)^{3/2}\\
&\times\exp\!\left(\frac{\mu_e+\mu_p-\mu_{H^0}}{T}\right)
\exp\!\left(-\frac{m_e+m_p-m_{H^0}}{T}\right).
\end{aligned}
\]
Chemical equilibrium sets the first exponential to unity. The mass
difference in the second is $m_e+m_p-m_{H^0}=E_I$, the hydrogen binding energy:
the bound atom is lighter than its separated constituents. Using
$m_{H^0}\simeq m_p$ \emph{only in the prefactor} gives
\[
\frac{n_en_p}{n_{H^0}}\simeq
\frac{g_eg_p}{g_{H^0}}\left(\frac{m_eT}{2\pi}\right)^{3/2}e^{-E_I/T}.
\]
The small binding-energy mass difference must be retained in the exponential.
For ground-state hydrogen, summing the spin states gives $g_e=2$, $g_p=2$,
$g_{H^0}=4$, so the degeneracy ratio is unity. Now explicitly substitute the
free and neutral densities in terms of the ionization fraction:
\[
\frac{n_en_p}{n_{H^0}}
=\frac{(x_en_H)(x_en_H)}{(1-x_e)n_H}
=n_H\frac{x_e^2}{1-x_e}.
\]
Equating the two expressions and dividing by $n_H$ gives the Saha equation
in this approximation,
\be
\boxed{
\frac{x_e^2}{1-x_e}=\frac{1}{n_H}
\left(\frac{m_eT}{2\pi}\right)^{3/2}e^{-E_I/T},
\qquad E_I=13.6\ {\rm eV}.}
\label{eq:Saha}
\ee
\paragraph{What sets the temperature?}
The guess $T\sim13.6$ eV comes from comparing a typical thermal photon energy with
the atomic ionization threshold. It is \emph{not} a conclusion of the Saha equation.
That equation also contains a large density-dependent prefactor. Define
$\eta_H=n_H/n_\gamma$, of order $10^{-9}$, and use the blackbody photon density:
\[
\frac{x_e^2}{1-x_e}=
\frac{\pi^2}{2\zeta(3)\eta_H}
\left(\frac{m_e}{2\pi T}\right)^{3/2}e^{-E_I/T}.
\]
At $T=E_I$ the right-hand side is enormous, implying $x_e\simeq1$. To make it of
order unity, the exponential must compensate the prefactor: one needs
$E_I/T\sim40$, not $E_I/T\sim1$. For example, setting $x_e=1/2$ and using a
representative cosmological hydrogen abundance gives an equilibrium temperature
near $0.32$ eV. Physically, there are roughly a billion photons per baryon: even
when a typical photon is too soft to ionize hydrogen, the high-energy tail can
keep the gas ionized.

The word \emph{recombination} refers to neutral atoms forming. The word \emph{decoupling}
refers to photons ceasing to scatter efficiently. The events are closely related but
not identical. Saha determines an \emph{equilibrium ionization fraction}, not the
last-scattering time. Accurate calculations also require the departure from Saha
equilibrium, excited atomic levels, the escape of Lyman-$\alpha$ photons, and
two-photon decays. Last scattering occurs near $z_*\simeq1100$, at
$T_*\simeq0.26$ eV (roughly $0.3$ eV). We use Saha for the equilibrium reasoning;
the scattering history below determines when photons can reach us.

\subsection{Optical depth and the visibility function}
Following the notation of the Particle Cosmology notes, define the optical depth from
conformal time $\tau$ to today by
\be
\tilde\tau(\tau)=\int_\tau^{\tau_0}a(\tau')n_e(\tau')\sigma_T\,d\tau'.
\label{eq:opticaldepth}
\ee
It follows from this definition that
\be
\dot{\tilde\tau}=-a n_e\sigma_T<0,
\qquad -\dot{\tilde\tau}=a n_e\sigma_T>0.
\ee
This $\tilde\tau$ is not conformal time; the tilde is important. The probability that
a photon travels from $\tau$ to us without scattering is $e^{-\tilde\tau}$. Therefore
\be
\boxed{g(\tau)=-\dot{\tilde\tau}\,e^{-\tilde\tau}}
\label{eq:visibility}
\ee
is the \emph{visibility function}: the probability density that a photon observed
today last scattered at time $\tau$. It is sharply peaked at $\tau_*$ but has a finite
width. Thus ``the last-scattering surface'' is a thin shell, not an infinitely thin
mathematical sphere.

\begin{center}
\begin{tabular}{@{}p{.29\linewidth}p{.63\linewidth}@{}}\toprule
Quantity & Approximate value\\ \midrule
Redshift of last scattering & $z_*\simeq1100$\\
Temperature & $T_*\simeq3000$ K $\simeq0.3$ eV\\
Age of the Universe & $t_*\simeq3.8\times10^5$ yr\\
Present mean CMB temperature & $T_0\simeq2.725$ K\\ \bottomrule
\end{tabular}
\end{center}

\keyresult{Recombination removes free electrons, the Thomson mean free path grows rapidly,
and the visibility function peaks. We see the CMB because the Universe changes from
opaque to transparent.}

% ============================================================
\section{From an exactly smooth Universe to anisotropies}

An exactly FLRW Universe would produce an exactly uniform CMB. We therefore perturb
the metric and matter fields around their homogeneous values,
\be
g_{\mu\nu}(\tau,\bx)&=&\bar g_{\mu\nu}(\tau)+\delta g_{\mu\nu}(\tau,\bx),\non
\rho_i(\tau,\bx)&=&\bar\rho_i(\tau)[1+\delta_i(\tau,\bx)].
\ee
The observed amplitude $10^{-5}$ motivates linear perturbation theory: keep only
terms first order in $\delta_i$, metric potentials, velocities, and $\Theta$.

Scalar, vector, and tensor perturbations evolve independently at linear order. Density
perturbations are scalar, so we focus on scalar modes. In conformal Newtonian gauge,
using precisely the convention of the Particle Cosmology notes,
\be
\boxed{ds^2=a^2(\tau)\left[-(1+2\Psi)d\tau^2
+(1-2\Phi)\delta_{ij}dx^idx^j\right].}
\label{eq:newtonian}
\ee
$\Psi$ acts like a Newtonian gravitational potential and $\Phi$ is the perturbation
of spatial curvature. Equation~\eqref{eq:newtonian} is a \emph{perturbed metric};
the perturbed Einstein equations determine how these potentials respond to matter.

\paragraph{Why just two scalar potentials?}
Before choosing coordinates, the scalar part of a metric perturbation contains four
functions: a time--time perturbation, a scalar time--space shift, a spatial trace,
and a scalar spatial shear. There are two scalar coordinate freedoms: a change of
time slicing and a longitudinal change of spatial coordinates. Newtonian gauge uses
them to set the shift and scalar shear to zero, leaving $\Psi$ and $\Phi$.
This counts scalar metric functions after a coordinate choice, \emph{not two
independently propagating gravitational degrees of freedom}. The scalar Einstein
constraint equations relate the potentials to density and momentum perturbations.
Matter supports scalar dynamics, including the sound waves we derive below. The two
freely propagating gravitational degrees of freedom of general relativity are the
two \emph{tensor} polarizations of gravitational waves, outside our present scalar
discussion.

\paragraph{What is anisotropic stress, and why do neutrinos produce it?}
Pressure is momentum flux: particles crossing a surface carry momentum through it.
An isotropic distribution has the same pressure in every direction. Its departure
from isotropy is the traceless part of the spatial stress tensor,
\[
T^i{}_j=P\delta^i{}_j+\pi^i{}_j,\qquad \pi^i{}_i=0.
\]
The tensor $\pi^i{}_j$ is called \emph{anisotropic stress}. In particle language it
comes from a quadrupole in the angular distribution of momenta, after removing the
isotropic pressure and bulk motion. A scalar perturbation can produce such a
quadrupole: its spatial gradient supplies a preferred direction locally, without
making the background Universe anisotropic.

After neutrinos decouple, they travel freely from regions with different densities
and temperatures. At a given point, particles arriving from different directions
therefore have different distributions. With few collisions to isotropize those
distributions, a quadrupole develops. Tightly coupled photons instead scatter often,
which suppresses their quadrupole; non-relativistic baryons and cold dark matter
have very small random-velocity stresses. Neutrinos are not unique in principle:
photons also develop anisotropic stress as tight coupling fails and after decoupling.
The traceless spatial Einstein equation relates this stress to $\Phi-\Psi$.
For negligible total scalar anisotropic stress, $\Phi=\Psi$; free-streaming neutrinos
make that equality approximate in the real Universe.

\subsection{Fourier modes simplify linear evolution}
A sound wave is a spatial pattern of alternating overdensities and underdensities.
We decompose an arbitrary small perturbation into sinusoidal patterns, or Fourier
modes. The vector $\bk$ specifies the direction of variation and its magnitude
$k=|\bk|$ the inverse length scale. Write, for example,
\be
\delta_\gamma(\tau,\bx)=\int\frac{d^3k}{(2\pi)^3}
\delta_\gamma(\tau,\bk)e^{i\bk\cdot\bx}.
\ee
Because the background is homogeneous and the equations are linear, different $\bk$
modes evolve independently. A spatial derivative becomes multiplication by $i\bk$;
coupled partial differential equations become ordinary differential equations in time
for each $k$.

The comoving wavelength is $\lambda_{\rm com}=2\pi/k$; $k^{-1}$ is the scale over
which the phase changes by one radian, not a full wavelength. Comparing the gradient
scale with the comoving Hubble length, we call a mode super-Hubble when $k\ll\Hs$
and sub-Hubble when $k\gg\Hs$. For acoustic evolution the more precise comparison
is with the sound horizon: appreciable phase evolution requires $kr_s\gtrsim1$.
At last scattering, a transverse scale of order $k^{-1}$ subtends an angle of order
$1/(k\chi_*)$. Lecture II will describe angular patterns by spherical-harmonic
multipoles $\ell$, with the approximate correspondence $\ell\sim k\chi_*$.

For photons $\rho_\gamma\propto T^4$. Linearizing gives
\be
\frac{\delta\rho_\gamma}{\bar\rho_\gamma}=4\frac{\delta T}{\bar T},
\qquad\boxed{\Theta_0=\frac14\delta_\gamma.}
\label{eq:ThetaMonopole}
\ee
Here $\Theta_0$ is the angle-averaged photon-temperature perturbation for a Fourier
mode. It should not be confused with today's mean temperature $T_0$.

More precisely, at each $(\tau,\bx)$ the photon perturbation also depends on the
propagation direction $\bn$. For a scalar Fourier mode define
$\mu=\hat{\bk}\cdot\bn$ and expand the angular dependence as
\be
\Theta(\tau,\bk,\mu)=\sum_{\ell=0}^{\infty}
(-i)^\ell(2\ell+1)\Theta_\ell(\tau,k)P_\ell(\mu).
\ee
The monopole $\Theta_0$ is the directional average and determines the photon energy
density. The dipole $\Theta_1$ is the forward--backward asymmetry and determines the
net photon flux, or momentum density. It is therefore often called the
\emph{photon temperature dipole}, the \emph{photon dipole}, or, more loosely, the
\emph{photon velocity}. With this multipole convention and the Fourier convention
used here,
\be
\boxed{\theta_\gamma=3k\Theta_1.}
\ee
The factor and sign are convention dependent. We reserve $\Theta_1$ for the angular
dipole, $v_\gamma$ for a scalar velocity amplitude, and $\theta_\gamma$ for the
velocity divergence. These are local moments of the photon momentum directions;
they are conceptually distinct from the observed whole-sky coefficients $a_{\ell m}$.

\subsection{Where did the perturbations come from?}
Here we take the initial conditions as given: an early mechanism such as
inflation produces an almost scale-invariant, nearly Gaussian, predominantly adiabatic
curvature perturbation $\Rs_i(\bk)$. ``Adiabatic'' means that all species initially
describe the same local time shift; there is a single growing scalar mode. On scales
larger than the Hubble radius, $\Rs$ is approximately conserved. Lecture II will use
its power spectrum $\Delta_{\Rs}^2(k)$, in the notation of the course notes.

Fourier modes make the three-dimensional evolution simple, even though the CMB is
observed on a sphere. Spherical harmonics provide the corresponding natural basis
for the two-dimensional observation. Lecture II connects the two descriptions.

% ============================================================
\section{The tightly coupled photon--baryon fluid}
\label{sec:acoustics}

\subsection{The physical picture}
Before recombination photons scatter so rapidly that their mean free path is much
smaller than the wavelength of interest. Photons and baryons share a common bulk
velocity,
\be
\bv_\gamma\simeq\bv_b.
\ee
Gravity tries to compress an overdense region. Photon pressure resists compression.
The result is the cosmological analogue of sound in air, except that photons provide
most of the pressure while photons and baryons both provide inertia.

The inertia of a relativistic fluid is set by its \emph{enthalpy density}, $\rho+P$:
to first order in velocity its momentum density is $(\rho+P)\bv$ in a local
orthonormal frame. Photons therefore contribute $4\bar\rho_\gamma/3$ to the inertia,
whereas non-relativistic baryons contribute approximately $\bar\rho_b$.
Define their ratio, the baryon-loading ratio used in the course notes,
\be
\boxed{R\equiv\frac{3\bar\rho_b}{4\bar\rho_\gamma}.}
\label{eq:R}
\ee
Thus \emph{baryon inertia} is the physical resistance of baryons to acceleration,
while \emph{baryon loading} describes how much of it is added relative to the photon
inertia. The total inertia is $(4\bar\rho_\gamma/3)(1+R)$.
Since $\rho_b\propto a^{-3}$ and $\rho_\gamma\propto a^{-4}$, $R\propto a$: the
baryon contribution becomes more important \emph{relative to photons}, even though
both densities decrease. In conformal time, $\dot R=\Hs R$.

For an adiabatic compression, the baryon-to-photon number ratio stays fixed. Since
$n_\gamma\propto T^3$ and $\rho_\gamma\propto T^4$, this gives
$\delta_b=3\delta_\gamma/4$, or $\delta\rho_b=R\,\delta\rho_\gamma$.
Using $\delta P_b\simeq0$ and $\delta P_\gamma=\delta\rho_\gamma/3$, the sound speed is
\be
c_s^2\equiv\frac{\delta P}{\delta\rho}
=\frac{\delta P_\gamma}{\delta\rho_\gamma+\delta\rho_b}
=\boxed{\frac{1}{3(1+R)}}.
\label{eq:soundspeed}
\ee
With no baryons, $c_s=1/\sqrt3$. Adding baryons lowers the sound speed because they add
inertia but essentially no pressure.

\subsection{Continuity and Euler equations}
We need two statements: how compression changes a density, and how pressure,
gravity, and collisions change a velocity. These are the energy and momentum
moments of the photon Boltzmann equation, or equivalently the continuity and Euler
equations of the fluid. Take one Fourier mode and define its velocity divergence
$\theta_i\equiv i\bk\cdot\bv_i$. Lowercase $\theta_i$ is a velocity-divergence variable;
uppercase $\Theta_0$ is a temperature perturbation. In real space, negative
divergence means converging flow and hence compression. A spatially uniform velocity
has zero divergence, so bulk translation by itself does not compress the fluid.

\paragraph{Continuity: compression changes photon temperature.}
In conformal Newtonian gauge the photon continuity equation is
\be
\dot\delta_\gamma=-\frac43\theta_\gamma+4\dot\Phi,
\label{eq:photoncontinuity}
\ee
The factor $4/3=(\bar\rho_\gamma+\bar P_\gamma)/\bar\rho_\gamma$ accounts for both
energy transport and pressure work. A useful way to see the metric term is that a
small coordinate cell has proper volume proportional to $a^3(1-3\Phi)$.
A changing $\Phi$ changes the local volume expansion; the effective perturbation
of its expansion rate is $\theta_\gamma-3\dot\Phi$. Thus the continuity equation
can be read as $\dot\delta_\gamma=-(4/3)(\theta_\gamma-3\dot\Phi)$.
Since $\Theta_0=\delta_\gamma/4$, it also says
\[
\dot\Theta_0=-\frac{\theta_\gamma}{3}+\dot\Phi.
\]
Compression raises the photon temperature, while changing spatial curvature also
changes the photon energy density measured locally.

\paragraph{Euler: what accelerates each component?}
The photon Euler equation, neglecting its quadrupole in tight coupling, is
\be
\dot\theta_\gamma=k^2\left(\frac14\delta_\gamma+\Psi\right)
+(-\dot{\tilde\tau})(\theta_b-\theta_\gamma).
\label{eq:photoneuler}
\ee
The pressure acceleration is $-\boldsymbol\nabla\delta P_\gamma/
(\bar\rho_\gamma+\bar P_\gamma)=-\boldsymbol\nabla\Theta_0$; the gravitational
acceleration is $-\boldsymbol\nabla\Psi$. Taking their divergences gives
$k^2\Theta_0$ and $k^2\Psi$, since $\nabla^2\to-k^2$. The remaining term is
Thomson drag. Its positive rate $-\dot{\tilde\tau}=a n_e\sigma_T$ is the number of
scatterings per unit \emph{conformal} time.

The baryon Euler equation, neglecting baryon thermal pressure, is
\be
\dot\theta_b+\Hs\theta_b=k^2\Psi
+\frac{-\dot{\tilde\tau}}{R}(\theta_\gamma-\theta_b).
\label{eq:baryoneuler}
\ee
Here $\Hs\theta_b$ is the expansion drag: in the absence of forces a
non-relativistic peculiar velocity decreases as $a^{-1}$. There is no corresponding
term for the photon bulk-velocity variable, whose radiation pressure and background
redshifting cancel it. Baryons feel gravity but contribute negligible pressure.

\paragraph{Why scattering equalizes velocities.}
If photons have a bulk drift relative to the electrons, their radiation is dipolar
in the electron rest frame. Scattering transfers momentum from this radiation flux
to electrons, and Coulomb interactions share it with the ions. Thus photons are
dragged toward the baryon velocity, and baryons toward the photon velocity.
The two \emph{force densities} are equal and opposite. Dividing them by the two
inertias gives different accelerations: the baryon acceleration contains
$(4\bar\rho_\gamma/3)/\bar\rho_b=1/R$. Indeed, the collision contribution to total
momentum obeys
\[
\frac{4\bar\rho_\gamma}{3}\dot\theta_\gamma\big|_{\rm coll}
+\bar\rho_b\dot\theta_b\big|_{\rm coll}=0.
\]
Opposite signs \emph{and this inertia weighting} are both necessary for momentum
conservation.

To check relaxation explicitly, define the velocity slip
$\Delta\theta=\theta_\gamma-\theta_b$. Retaining only the collision terms gives
\[
\left.\frac{d\Delta\theta}{d\tau}\right|_{\rm coll}
=-(-\dot{\tilde\tau})\left(1+\frac1R\right)\Delta\theta.
\]
The coefficient on the right is negative: either sign of slip decays exponentially
when the rate is approximately constant. Pressure and expansion continually source
a small slip, but rapid scattering keeps it small. Tight coupling requires the
scattering rate to exceed both the expansion and mode-evolution rates (in
particular $-\dot{\tilde\tau}\gg\Hs,k$ is sufficient).
It locks \emph{bulk velocities}, not the density contrasts or the temperature
perturbations of the two species.

The original Particle Cosmology notes instead write the scalar velocity as
\be
\bv_i=i\hat{\bk}\,v_i.
\ee
With the Fourier convention used here, the translation is
\be
\boxed{\theta_i=i\bk\cdot\bv_i=-k v_i.}
\ee
Thus $v_i$ in the long notes and $\theta_i$ here contain the same scalar velocity
information; the extra factor and sign come only from the definition.

\subsection{Deriving the acoustic oscillator}
\paragraph{First derive the equation for the actual photon temperature.}
Add Eq.~\eqref{eq:photoneuler} to $R$ times Eq.~\eqref{eq:baryoneuler}.
The collision terms cancel. This order matters: a small velocity slip multiplied
by a large scattering rate can transfer a finite momentum. Setting the collision
terms to zero in each species equation separately would lose that transfer.
Only after adding the equations take the leading tight-coupling limit
$\theta_\gamma\simeq\theta_b\equiv\theta$ to obtain
\[
(1+R)\dot\theta+\Hs R\theta
=k^2\Theta_0+(1+R)k^2\Psi.
\]
Gravity accelerates the whole fluid; only photons supply the pressure force.
Dividing by $1+R$ makes the reduced pressure acceleration explicit:
\[
\dot\theta+\frac{\dot R}{1+R}\theta
=\frac{k^2}{1+R}\Theta_0+k^2\Psi,
\qquad \dot R=\Hs R.
\]
Now differentiate $\dot\Theta_0=-\theta/3+\dot\Phi$, substitute the combined Euler
equation, and use $\theta=-3(\dot\Theta_0-\dot\Phi)$. This gives
\[
\boxed{\ddot\Theta_0+\frac{\dot R}{1+R}\dot\Theta_0+k^2c_s^2\Theta_0
=\ddot\Phi+\frac{\dot R}{1+R}\dot\Phi-\frac{k^2}{3}\Psi.}
\]
This equation is already a complete temperature oscillator at leading order in
tight coupling, for specified gravitational potentials. The pressure term restores
the fluid, $\dot R/(1+R)$ accounts for the changing inertia, $-k^2\Psi/3$ is
gravitational compression, and the derivatives of $\Phi$ arise from the changing
spatial metric in the continuity equation.

\paragraph{Why introduce an effective temperature?}
The intrinsic perturbation $\Theta_0$ tells us how hot the plasma is \emph{where the
photon is emitted}. The photon also undergoes a gravitational redshift on its way
out. For time-independent potentials, after factoring out the background redshift,
the fractional energy change between emission and observation is
$\Psi_*-\Psi_{\rm obs}$. The observer term is common to all directions and contributes
only a monopole. The emission term therefore adds $\Psi_*$ to the intrinsic
temperature perturbation. This motivates
\be
X\equiv\Theta_0+\Psi.
\ee
In our sign convention a potential well has $\Psi<0$, so its redshift reduces the
temperature received from an intrinsically hot region. Thus $X$ is the intrinsic
temperature \emph{plus the local gravitational contribution}, not a new physical
degree of freedom. It is the emission contribution used in Lecture II; Doppler
shifts and evolving potentials along the path add further terms to the observable.
We could keep using $\Theta_0$, but $X$ makes both the observed combination and the
gravity-shifted oscillation particularly transparent.

Substitute $\Theta_0=X-\Psi$ into the temperature equation:
\be
\boxed{
\begin{aligned}
\ddot X+\frac{\dot R}{1+R}\dot X+k^2c_s^2X
={}&\ddot\Phi+\ddot\Psi
+\frac{\dot R}{1+R}(\dot\Phi+\dot\Psi)\\
&{}-k^2c_s^2R\Psi.
\end{aligned}}
\label{eq:oscillator}
\ee
This is a damped, driven harmonic oscillator for every Fourier mode $k$:
\begin{itemize}
\item $kc_s$ is the acoustic frequency in conformal time; the restoring force
originates in photon pressure;
\item $\dot R/(1+R)$ changes the amplitude as the baryon inertia grows. This is
expansion-related damping, not dissipative photon diffusion;
\item $-k^2c_s^2R\Psi$ shifts the equilibrium even for constant potentials;
\item time derivatives of $\Phi$ and $\Psi$ drive the oscillation when gravity evolves.
\end{itemize}

\paragraph{A solvable limit and the acoustic phase.}
For intuition, freeze $R$, $\Phi$, and $\Psi$ over an oscillation, so their time
derivatives vanish. Equation~\eqref{eq:oscillator} reduces to
\[
\ddot X+k^2c_s^2(X+R\Psi)=0.
\]
The displacement from equilibrium is therefore $X+R\Psi$, whose solution is
\be
X(\tau)+R\Psi=C_1\cos[k r_s(\tau)]+C_2\sin[k r_s(\tau)],
\label{eq:oscillatorsolution}
\ee
where
\be
\boxed{r_s(\tau)=\int_0^\tau c_s(\tau')d\tau'}
\label{eq:soundhorizon}
\ee
is the comoving sound horizon: the maximum comoving distance a sound wave could have
travelled by time $\tau$. For the constant-coefficient toy model,
$r_s=c_s\tau$ (with the chosen time origin), and
Eq.~\eqref{eq:oscillatorsolution} is exact. When $R$ varies slowly, the accumulated
phase is still approximately $k\int c_s d\tau=kr_s$, but the amplitude and the
equilibrium also evolve: constant $C_1,C_2$ are then only a toy approximation.
Rapid potential evolution requires solving the driven equation.

For the dominant adiabatic growing mode, all wavelengths begin with a coherent phase:
they are initially displaced with negligible velocity. Consequently the cosine term
dominates. At recombination we take a snapshot. Modes satisfying
\be
k_n r_s(\tau_*)\simeq n\pi,
\qquad n=1,2,3,\ldots
\label{eq:peakcondition}
\ee
are at extrema of compression or rarefaction. This harmonic series in $k$ becomes the
series of acoustic peaks in angular multipole $\ell$ after projection onto the sky.

\keyresult{Gravity displaces the photon--baryon fluid, photon pressure restores it,
inflationary/adiabatic initial conditions synchronize the phase, and recombination
takes a snapshot. These four facts produce the acoustic peaks.}

\subsection{Baryon loading and gravitational driving}
We can now derive the equilibrium shift rather than assume it. For a static
potential and fixed $R$, set the velocity and acceleration to zero in the combined
Euler equation. Pressure must balance gravity:
\[
0=\frac{k^2}{1+R}\Theta_{0,\rm eq}+k^2\Psi,
\qquad
\Theta_{0,\rm eq}=-(1+R)\Psi.
\]
Photons must provide enough pressure to support \emph{both} their own inertia and
the baryons in the potential well. More baryons therefore require a larger intrinsic
temperature and density perturbation for balance. Converting to the effective
temperature gives
\[
\boxed{X_{\rm eq}=\Theta_{0,\rm eq}+\Psi=-R\Psi.}
\]
With no baryons the intrinsic equilibrium temperature $-\Psi$ is exactly cancelled
by the gravitational redshift, so $X_{\rm eq}=0$. Baryons leave the extra offset
$-R\Psi$. For a potential well ($\Psi<0$), this offset is positive. An oscillation
of amplitude $A$ has extrema $X=-R\Psi\pm A$: its compressed, hot phase is enhanced
relative to its rarefied, cold phase. Projecting these coherent oscillations gives
the characteristic odd--even temperature-peak modulation: odd compression peaks
are enhanced as the baryon density increases.

During exact matter domination, the gravitational potentials are nearly constant.
Around radiation--matter equality they evolve. A time-dependent equilibrium point
does work on the oscillator. Modes entering the horizon during radiation domination
are therefore boosted by \emph{radiation driving}. This is why the peak envelope also
contains information about the physical matter density and the time of equality.

The acoustic peaks are evidence for coherent initial conditions. If different
regions or Fourier modes began oscillating with random phases, averaging over them
would wash out the peak--trough sequence.

% ============================================================
\section{Why the peaks do not continue forever}

Tight coupling is an approximation. Photons have a small but nonzero mean free path
and random-walk out of hot overdense regions. Heat conduction and shear viscosity erase
anisotropy below a diffusion length. In Fourier space the acoustic amplitude is
approximately multiplied by
\be
\exp[-k^2/k_D^2].
\ee
This is Silk, or photon-diffusion, damping. The finite thickness of the visibility
function adds further averaging: photons last-scattering at slightly different times
sample slightly different oscillation phases.

The main physical scales at last scattering are therefore
\begin{itemize}
\item the comoving distance $\chi_*$ from us to the last-scattering surface;
\item the sound horizon $r_s(\tau_*)$, setting the acoustic spacing;
\item the diffusion length $r_D\sim k_D^{-1}$, setting the damping tail;
\item the horizon around equality, setting gravitational driving.
\end{itemize}
Projection turns the sound horizon into an angle,
\be
\boxed{\theta_s=\frac{r_s(\tau_*)}{\chi_*},
\qquad \ell_A\simeq\frac{\pi}{\theta_s}.}
\ee
This compact equation is why the CMB constrains geometry: early-time physics predicts
the ruler $r_s$, while late-time expansion determines the distance $\chi_*$.

\section*{Summary}
We can now tell the CMB story without equations:
\begin{enumerate}
\item An FLRW Universe expands and cools; photon wavelengths grow as $a$ and $T\propto a^{-1}$.
\item The early Universe is an opaque photon--baryon plasma maintained by Thomson scattering.
\item Neutral hydrogen formation rapidly lowers $n_e$; the visibility function peaks and photons free-stream.
\item Small primordial perturbations provide gravitational potential wells.
\item Gravity and photon pressure make the tightly coupled plasma oscillate with sound speed
$c_s^2=1/[3(1+R)]$.
\item Recombination photographs these coherently phased oscillations. Baryons modulate
compressions and diffusion damps short wavelengths.
\item We observe the three-dimensional pattern projected onto a two-dimensional sphere.
Lecture II turns this statement into the angular spectrum $C_\ell$.
\end{enumerate}

\section*{Companion numerical notebook}
The physical sequence in this lecture can be explored in the accompanying notebook:
\begin{center}
\href{https://colab.research.google.com/github/daanmeerburg/RTG_3186_CMB_Lectures/blob/main/CMB_physics_from_expansion_to_acoustics.ipynb}{\textbf{Open the Lecture I exercise directly in Google Colab}}.
\end{center}
It visualizes radiation, matter, and vacuum domination; solves the Saha equation;
constructs an approximate visibility function; calculates the photon--baryon sound
horizon; and plots the coherent density and velocity phases. The notebook emphasizes
where the Saha approximation differs from a precision recombination calculation and
ends with exercises exploring the physical parameters.

\section*{Further derivation A: scale-factor solutions}
For a component with constant $w$, combine $\rho\propto a^{-3(1+w)}$ with Friedmann:
\[
\frac{1}{a}\frac{da}{dt}\propto a^{-3(1+w)/2}.
\]
For $w\ne-1$, separation of variables gives
\be
a(t)\propto t^{\frac{2}{3(1+w)}}.
\ee
Thus radiation gives $a\propto t^{1/2}$ and matter gives $a\propto t^{2/3}$.

\section*{Further derivation B: why the large-scale temperature is gravitational}
The photon reaching us carries an intrinsic temperature perturbation and is redshifted
while climbing out of a potential well. At last scattering these combine as
\be
\left[\frac14\delta_\gamma+\Psi\right]_*.
\ee
For an adiabatic mode on scales much larger than the sound horizon during matter
domination, the result is the ordinary Sachs--Wolfe relation
\be
\left.\frac{\delta T}{T}\right|_{\rm SW}\simeq\frac13\Psi_*.
\ee
There is also a Doppler contribution from the baryon velocity and an integrated
Sachs--Wolfe contribution if $\Phi$ or $\Psi$ evolves along the photon path. These are
the opening ingredients of Lecture II.

\section*{Notation guide}
\begin{center}
\begin{tabular}{@{}p{.33\linewidth}p{.62\linewidth}@{}}\toprule
Symbol & Meaning\\ \midrule
$t$, $H=a^{-1}da/dt$ & cosmic time and cosmic-time Hubble rate; dot means $d/dt$ in Sec.~2\\
$\tau$, $\Hs=\dot a/a=aH$ & conformal time and conformal Hubble rate; thereafter dot means $d/d\tau$\\
$\tilde\tau$ & Thomson optical depth, not a time coordinate\\
$\chi$, $\chi_*$ & comoving radial coordinate and distance to last scattering\\
$n_H$, $n_{H^0}$, $x_e$ & total hydrogen-nucleus density, neutral-atom density, and $n_e/n_H$\\
$\Phi,\Psi$ & conformal-Newtonian-gauge metric potentials\\
$\Theta=\delta T/\bar T$ & fractional photon-temperature perturbation\\
$\delta_\gamma$ & fractional photon-energy-density perturbation; $\Theta_0=\delta_\gamma/4$\\
$\Theta_1$ & photon temperature dipole; proportional to photon momentum or bulk flow\\
$\theta_i=i\bk\cdot\bv_i$ & velocity divergence, distinct from temperature $\Theta_0$\\
$X=\Theta_0+\Psi$ & intrinsic temperature plus local gravitational contribution\\
$R=3\bar\rho_b/(4\bar\rho_\gamma)$ & baryon-to-photon inertia ratio (baryon loading)\\
$\Rs$, $\Delta_{\Rs}^2$ & comoving curvature perturbation and its dimensionless power spectrum\\
$r_s$ & comoving sound horizon\\ \bottomrule
\end{tabular}
\end{center}

\section*{Source and scope note}
This lecture is a shortened, reordered version of the Particle Cosmology course notes:
FLRW dynamics and redshift from Lectures 1--2; equilibrium, decoupling, and recombination
from Lecture 3; scalar perturbations and conformal Newtonian gauge from Lecture 4; and
tight coupling and CMB acoustic physics from Lecture 5. The notation has been aligned
with those notes. Material has been reordered so that no prior cosmology course is
assumed; detailed gauge transformations and the full Boltzmann hierarchy are postponed.
The intermediate steps on chemical equilibrium, momentum exchange, and the
oscillator's equilibrium have been expanded here for a first introduction.
For supplementary derivations see
\href{https://arxiv.org/abs/astro-ph/9506072}{Ma \& Bertschinger (1995)} for the
fluid equations and
\href{https://background.uchicago.edu/~whu/araa/node10.html}{Hu \& Dodelson's
discussion of baryon loading} for the oscillator interpretation. Metric signs
throughout follow the Particle Cosmology master notes.

\end{document}
